\clearpage

\section{Phase Diagram}\label{sec-nbcp-phases}

The classical three-sublattice analysis supplies reference branches and
local boundaries. A phase diagram requires comparison with other
magnetic cells and the relevant quantum energies. The distinction is
maintained below: a stable magnon spectrum is a local test, while the
lowest energy among an incomplete set of candidate states need not
identify the global ground state.

\subsection{Classical phase diagram}\label{sec-nbcp-classical-phases}

\subsubsection{Three-sublattice classical
energy}\label{three-sublattice-classical-energy}

Let \(\mathbf n_\alpha=\mathbf S_\alpha/S\), with \(\alpha=A,B,C\),
specify a uniform three-sublattice configuration. Each pair of
sublattices is connected by all three bond orientations. Since
\(\sum_d\mathsf J_d=3\operatorname{diag}(J,J,J_z)\), its classical
energy per physical spin reduces to

\begin{equation}\phantomsection\label{eq-nbcp-three-energy}{
e_{\rm cl}=S^2\sum_{(\alpha\beta)=(AB),(BC),(CA)}
\left[J\sin\vartheta_\alpha\sin\vartheta_\beta
\cos(\phi_\alpha-\phi_\beta)
+J_z\cos\vartheta_\alpha\cos\vartheta_\beta\right]
-\frac{hS}{3}\sum_\alpha\cos\vartheta_\alpha.
}\end{equation}

Here \(\vartheta\) may be signed, with
\(\mathbf n=(\sin\vartheta\cos\phi,\sin\vartheta\sin\phi,\cos\vartheta)\).
Negative \(\vartheta\) is a coordinate choice, equivalent to a positive
polar angle and an azimuth shifted by \(\pi\).

The SOC contributions cancel in this restricted classical energy. They
do not disappear from the operator Hamiltonian, from finite-momentum
fluctuations, or from states with other magnetic unit cells. In
particular, the cancellation does not establish SOC-independent global
phase boundaries.

\subsubsection{Classical reference fields and their stability
meaning}\label{classical-reference-fields-and-their-stability-meaning}

The UUD stability matrix can be derived in regular Cartesian coordinates
at the collinear poles. Put
\(\mathbf S_A=(\sqrt S x_A,\sqrt S y_A,S-(x_A^2+y_A^2)/2)\) and
similarly for B; for C replace the last component by
\(-S+(x_C^2+y_C^2)/2\). To quadratic order,
\(\delta e=(\mathbf x^{\mathsf T}M\mathbf x+\mathbf y^{\mathsf T}M\mathbf y)/2\),
where

\begin{equation}\phantomsection\label{eq-nbcp-uud-hessian}{
M=\begin{pmatrix}
h/3&JS&JS\\
JS&h/3&JS\\
JS&JS&2J_zS-h/3
\end{pmatrix}.
}\end{equation}

The antisymmetric A/B eigenvector has eigenvalue \(h/3-JS\). The
symmetric A/B sector and C reduce to a \(2\times2\) matrix with diagonal
entries \(h/3+JS\) and \(2J_zS-h/3\), and off-diagonal entry
\(\sqrt2JS\). Setting the lower eigenvalue to zero gives the upper UUD
boundary. At P, the diagonal entries become \(h/3-2J_zS\), giving the
saturation boundary. For \(J_z>J>0\) the three reference fields are

\begin{equation}\phantomsection\label{eq-nbcp-critical-fields}{
\begin{aligned}
h_{c1}&=3SJ,\\
h_{c2}&=3S\left[J_z-\frac J2+
\sqrt{J_z^2+J_zJ-\frac74J^2}\right],\\
h_{c3}&=3S(J+2J_z),\qquad B_{ci}=h_{ci}/(g_z\mu_B).
\end{aligned}
}\end{equation}

These are the classical XXZ reference boundaries and uniform
three-sublattice softening conditions. Quantum corrections, other
ordering wave vectors and SOC-driven competitors require separate
calculations.

\begin{longtable}[]{@{}
  >{\raggedleft\arraybackslash}p{(\linewidth - 4\tabcolsep) * \real{0.3333}}
  >{\raggedright\arraybackslash}p{(\linewidth - 4\tabcolsep) * \real{0.3333}}
  >{\raggedright\arraybackslash}p{(\linewidth - 4\tabcolsep) * \real{0.3333}}@{}}
\toprule\noalign{}
\begin{minipage}[b]{\linewidth}\raggedleft
\(J\) (meV)
\end{minipage} & \begin{minipage}[b]{\linewidth}\raggedright
\((h_{c1},h_{c2},h_{c3})\) (meV)
\end{minipage} & \begin{minipage}[b]{\linewidth}\raggedright
\((B_{c1},B_{c2},B_{c3})\) (T)
\end{minipage} \\
\midrule\noalign{}
\endfirsthead
\toprule\noalign{}
\begin{minipage}[b]{\linewidth}\raggedleft
\(J\) (meV)
\end{minipage} & \begin{minipage}[b]{\linewidth}\raggedright
\((h_{c1},h_{c2},h_{c3})\) (meV)
\end{minipage} & \begin{minipage}[b]{\linewidth}\raggedright
\((B_{c1},B_{c2},B_{c3})\) (T)
\end{minipage} \\
\midrule\noalign{}
\endhead
\bottomrule\noalign{}
\caption{Classical three-sublattice reference fields for the stated exchange parameters.}\label{tab-nbcp-03-phase-diagram-1}\\
\endlastfoot
0.075 & (0.112500, 0.315916, 0.487500) & (0.418417, 1.174976,
1.813142) \\
0.076 & (0.114000, 0.314316, 0.489000) & (0.423996, 1.169024,
1.818721) \\
\end{longtable}

Both rows use \(J_z=0.125\) meV, \(S=1/2\), \(g_z=4.645\) and
\(\mu_B=0.05788381806\) meV/T. The older notes also quote
\((0.1005,0.3278,0.4755)\) meV while specifying \(J=0.076\) meV. Those
numbers are not evaluations of Eq.~\ref{eq-nbcp-critical-fields} at the
stated inputs and are not used here. The present reference fields were
recalculated; the archived phase scans were not rerun.

\subsection{Quantum corrections: zero-point
energy}\label{sec-nbcp-phase-quantum-correction}

At harmonic order, the comparison adds the spin-wave vacuum energy to
the classical energy for each admissible reference state. The positive
physical magnon bands and the normal-ordering trace constant must both
be included, and all candidate energies must use the same per-spin
normalization and Brillouin-zone convention. An unstable quadratic
background cannot be assigned a physical harmonic vacuum energy by
discarding its unstable modes.

The existing calculations resolve the vacuum-energy dependence along
selected Y/V angular orbits; they do not reoptimize and compare all
candidate phases across the model's parameter space. Consequently, the
classical reference fields are not promoted to quantum-corrected phase
boundaries. The local angular selection and its pseudo-Goldstone gap are
treated in Section~\ref{sec-nbcp-yv-low-energy}; numerical quadrature
and convergence are recorded in Section~\ref{sec-convergence}.

\subsection{Characterization of
phases}\label{sec-nbcp-phase-characterization}

\subsubsection{Y, UUD, V and the polarized
state}\label{y-uud-v-and-the-polarized-state}

At common azimuth \(\phi=0\), the following signed-angle definitions fix
the states independently of historical names:

\begin{longtable}[]{@{}
  >{\raggedright\arraybackslash}p{(\linewidth - 4\tabcolsep) * \real{0.3333}}
  >{\raggedright\arraybackslash}p{(\linewidth - 4\tabcolsep) * \real{0.3333}}
  >{\raggedright\arraybackslash}p{(\linewidth - 4\tabcolsep) * \real{0.3333}}@{}}
\toprule\noalign{}
\begin{minipage}[b]{\linewidth}\raggedright
State
\end{minipage} & \begin{minipage}[b]{\linewidth}\raggedright
\((\vartheta_A,\vartheta_B,\vartheta_C)\)
\end{minipage} & \begin{minipage}[b]{\linewidth}\raggedright
Characteristic structure
\end{minipage} \\
\midrule\noalign{}
\endfirsthead
\toprule\noalign{}
\begin{minipage}[b]{\linewidth}\raggedright
State
\end{minipage} & \begin{minipage}[b]{\linewidth}\raggedright
\((\vartheta_A,\vartheta_B,\vartheta_C)\)
\end{minipage} & \begin{minipage}[b]{\linewidth}\raggedright
Characteristic structure
\end{minipage} \\
\midrule\noalign{}
\endhead
\bottomrule\noalign{}
\caption{Signed-angle definitions of the three-sublattice reference states.}\label{tab-nbcp-03-phase-diagram-2}\\
\endlastfoot
Y & \((t,-t,\pi)\) & Opposite transverse components; one spin
antiparallel to the field \\
UUD & \((0,0,\pi)\) & Two up spins and one down spin; \(m_z=S/3\) \\
V & \((t,t,v)\) & Two equal canted spins; the third has the opposite
transverse direction \\
P & \((0,0,0)\) & All spins parallel to the field; \(m_z=S\) \\
\end{longtable}

For Y, substituting \(c=\cos t\) into Eq.~\ref{eq-nbcp-three-energy}
gives

\begin{equation}\phantomsection\label{eq-nbcp-y-energy}{
e_Y=S^2[-J+(J+J_z)c^2-2J_zc]-\frac{hS}{3}(2c-1),
\qquad c=\frac{h+3SJ_z}{3S(J+J_z)}.
}\end{equation}

Its merger with UUD occurs at \(c=1\). The collinear energies are
\(e_{\rm UUD}=-S^2J_z-hS/3\) and \(e_P=3S^2J_z-hS\). Their crossing
alone does not locate the intervening canted phase.

For the V family, the energy and two independent stationarity conditions
are

\begin{equation}\phantomsection\label{eq-nbcp-v-stationarity}{
\begin{aligned}
e_V={}&S^2\left[J(\sin^2t+2\sin t\sin v)
+J_z(\cos^2t+2\cos t\cos v)\right]
-\frac{hS}{3}(2\cos t+\cos v),\\
0={}&J\cos t(\sin t+\sin v)-J_z\sin t(\cos t+\cos v)
+\frac{h}{3S}\sin t,\\
0={}&2J\sin t\cos v-2J_z\cos t\sin v+\frac{h}{3S}\sin v.
\end{aligned}
}\end{equation}

The second and third lines are respectively
\((\partial_t e_V)/(2S^2)=0\) and \((\partial_v e_V)/S^2=0\). Their
solutions must also pass a stability check. At \(B=1.4\) T with the
current gap parameters, the selected classical branch has
\((t,t,v)=(0.409364861,0.409364861,-1.223629194)\) rad.

The original {\small\texttt{\nolinkurl{Three_MSL.tex}}} calls a family
with two equal canted spins a ``Gamma phase.'' Its normalized vectors
are

\[
\mathbf n_A=\mathbf n_B=(-\sin t,0,\cos t),
\qquad \mathbf n_C=(\sin\psi,0,-\cos\psi).
\]

A common \(R_z(\pi)\) rotation maps this family to V above with
\(v=\psi-\pi\). This is a geometric identification of the reference
families; at nonzero \(J_\Gamma\), that spin-only rotation is not a
symmetry of the full Hamiltonian. The old ``V'' ansatz instead assigns
opposite transverse components to A and B. Its restricted extrema must
not be substituted for the V solution above or identified with the
paper's \(\Psi\) phase merely by renaming.

\subsubsection{Competing magnetic cells and quantum
results}\label{competing-magnetic-cells-and-quantum-results}

The source notebook contains a two-sublattice stripe energy and a
four-sublattice variational construction. These are essential
competitors because they retain the anisotropic bond contributions that
cancel in the three-sublattice energy. For a stripe whose
same-sublattice bonds are of type 0, the bond count gives

\begin{equation}\phantomsection\label{eq-nbcp-stripe-energy}{
E^{(2)}_{\rm cell}=\mathbf S_A^{\mathsf T}\mathsf J_0\mathbf S_A
+\mathbf S_B^{\mathsf T}\mathsf J_0\mathbf S_B
+2\mathbf S_A^{\mathsf T}(\mathsf J_1+\mathsf J_2)\mathbf S_B
-h(S_A^z+S_B^z).
}\end{equation}

Comparison with a three-sublattice state requires
\(E^{(2)}_{\rm cell}/2\) and \(E^{(3)}_{\rm cell}/3\). A stationary
point found in a fixed spin plane still requires fluctuations outside
that plane and at finite momentum to be checked. In particular, an
exactly in-plane stripe at \(h\ne0\) experiences a linear Zeeman torque;
a field-independent Hessian evaluated there does not prove finite-field
stability.

The target paper reports Y, UUD, V, \(\Psi\) and P at zero
bond-dependent SOC, and stripe and four-site SkX orders as SOC grows.
Its phase map uses iDMRG cylinders, while the reported \(\Psi\) state is
not a stable LSWT reference. These results motivate the competitor list;
their phase boundaries and finite-circumference convergence have not
been independently reproduced here. See
Park et al., Section II~\cite{park2026}.
